\section{Admission of the candidate force to the forced theorem}
\label{r:sec:force-extension}

Part~II applies after the momentum residual of the proposed field has been
placed in its prescribed force class. The cutoff calculation below keeps the
full localization error, proves uniform bounds for every fixed mixed
derivative, and produces compatible compactly supported terminal jets.

\subsection{The localization residual}

Multiplication by a spatial cutoff changes the fluid's velocity gradients.
The resulting force must include those changes, including the correction
needed to preserve incompressibility. OpenAI's local representation is
\(v=\operatorname{curl}\mathcal A+B e_\theta\), with \(B\) independent
of the angle. Its localization is exactly
\begin{equation}
 u=cv+w,\qquad w=\nabla c\times\mathcal A,\qquad p=c\pi,
 \qquad c=\chi_x\chi_t.
 \label{r:eq:localization-fields}
\end{equation}
Here \(\pi\) is the local pressure, and \(c\) is axisymmetric.
This is \cite[(10.4)]{openai-paper}: the curl acts on \(c\mathcal A\),
while the swirl is multiplied directly by \(c\). In particular,
\(\nabla\cdot u=0\). Introducing a vector potential for the exterior
swirl and localizing that additional potential would change the field.

Write \(\mathcal R_\nu(v,\pi)=v_t+(v\cdot\nabla)v-\nu\Delta v+\nabla\pi\).
Expanding \eqref{r:eq:localization-fields} gives the entire localization error:
\begin{equation}
\begin{aligned}
 f-c\mathcal R_\nu(v,\pi)
 ={}&(c_t-\nu\Delta c)v-2\nu\sum_{i=1}^3(\partial_i c)\partial_i v
       +\pi\nabla c\\
 &+(c^2-c)(v\cdot\nabla)v+c(v\cdot\nabla c)v\\
 &+w_t-\nu\Delta w+c(v\cdot\nabla)w\\
 &+c(w\cdot\nabla)v+(w\cdot\nabla)w.
\end{aligned}
\label{r:eq:localization-residual}
\end{equation}
The apparent additional product \((w\cdot\nabla c)v\) vanishes because
\((\nabla c\times\mathcal A)\cdot\nabla c=0\).
Every remaining term belongs to the prescribed force. Where
both cutoffs are constant at one, \(w=0\) and the error vanishes.
Where they vanish on a neighborhood, the localized fields vanish too.
The extra force is confined to the transitions between these regions.

In the heat exterior, near time one, \(\mathcal A=0\) and
\(v=K e_\theta\). The formula becomes a concrete force field:
\begin{equation}
\begin{aligned}
 f={}&\left[\chi_x(1-\chi_x)\frac{K^2}{r}
                 +\pi\partial_r\chi_x\right]e_r
        +\pi\partial_z\chi_x e_z\\
 &-\nu\left[2(\partial_r\chi_x)K_r
       +K\left(\partial_{rr}\chi_x+r^{-1}\partial_r\chi_x
                         +\partial_{zz}\chi_x\right)\right]e_\theta.
\end{aligned}
\label{r:eq:exterior-cutoff-force-full}
\end{equation}
The radial force accounts for the changed centrifugal balance and the
pressure cutoff. The axial term comes from the pressure cutoff alone.
The angular force accounts for diffusion across the tapered swirl.
All three vanish where \(\chi_x\) is constant at one. The bounds below
control every finite derivative on their transition support.

For a space--time region \(\Sigma\), define the finite mixed-jet
seminorm
\[
 \mathcal J_m(F;\Sigma)
 =\sum_{|\gamma|\le m}\frac{\norm{\partial^\gamma F}_{L^\infty(\Sigma)}}{\gamma!},
 \qquad \gamma\in\mathbb N_0^4.
\]
The four entries of \(\gamma\) count three spatial derivatives and one
temporal derivative. The multinomial product rule gives
\(\mathcal J_m(FG)\le\mathcal J_m(F)\mathcal J_m(G)\), also for dot
and cross products. Each coordinate derivative obeys
\(\mathcal J_m(\partial_i F)\le(m+1)\mathcal J_{m+1}(F)\).

On a transition region put
\[
 \begin{aligned}
 C_m&=\mathcal J_m(c),& V_m&=\mathcal J_m(v),\\
 A_m&=\mathcal J_m(\mathcal A),& P_m&=\mathcal J_m(\pi).
 \end{aligned}
 \qquad
 U_m=C_mV_m+3(m+1)C_{m+1}A_m.
\]
Then \(\mathcal J_m(u)\le U_m\) and
\(\mathcal J_m(p)\le C_mP_m\). Differentiating the full momentum
residual gives the explicit force-jet bound
\begin{equation}
\begin{aligned}
 \mathcal J_m(f)\le{}&(m+1)U_{m+1}
       +3(m+1)U_mU_{m+1}\\
 &+3\nu(m+1)(m+2)U_{m+2}
       +3(m+1)C_{m+1}P_{m+1}.
\end{aligned}
\label{r:eq:localized-jet-bound}
\end{equation}
This bounds every fixed temporal and spatial source derivative once the
local field jets on the transition have been bounded. It requires no
constant shared by infinitely many orders.

\subsection{Uniform derivative bounds on the cutoff transitions}

The relevant bounds follow from the actual support geometry. The cutoff
in \cite[p.~118]{openai-paper} equals one for
\(r\le r_0/2\), \(|z|\le z_0/2\), and has compact support inside
\(r<r_0\), \(|z|<z_0\). Every spatial cutoff error in
\eqref{r:eq:localization-residual}, including \(c^2-c\) near time one,
is supported where
\[
 r\ge r_0/2\quad\hbox{or}\quad |z|\ge z_0/2.
\]
The construction's coordinates satisfy
\[
 \tau=q(1-\eta^2),\quad z=q^D\eta,\quad
 X=\frac{r^2}{2q},\quad |\eta|\le1,
 \qquad D=\frac12-h>0.
\]
On the axial transition, \(q\ge |z|^{1/D}\) yields
\[
 q\ge q_z:=(z_0/2)^{1/D}>0.
\]
This remains true as \(r\) tends to zero. On the radial transition,
any correction inside \(X\le X_{\rm ext}\) satisfies
\[
 q=\frac{r^2}{2X}\ge q_r:=\frac{r_0^2}{8X_{\rm ext}}>0.
\]
Thus every interior correction meeting a spatial transition has
\(q\ge\min(q_z,q_r)>0\). The cutoff choices also ensure
\(q<q_*/2\), by \cite[(10.3)]{openai-paper}. This is precisely the
closed scale range covered by Theorem~3.1(ii) of that paper.

The one-sided interior bound has a concrete finite-order meaning.
On \(q\ge\delta>0\), only finitely many summands in
\cite[(9.21)]{openai-paper} survive, since their cutoff scales tend to
infinity. The finite representatives, fixed phase labels, and closed
parameter bounds are those specified in Proposition~9.9, Step~3.
Coordinate conversion has denominator
\(1-2h\eta^2\ge1-2h>0\). At the axis, the potential is represented
in Cartesian coordinates as \(a(r^2,z,t)(-x_2,x_1,0)\), as in
\cite[(10.1)]{openai-paper}. These are the stated interior inputs to
the finite numbers \(V_m,A_m,P_m\) in \eqref{r:eq:localized-jet-bound}.

If \(q\) approaches zero on the remaining radial transition, the field
is instead in \(X\ge X_{\rm ext}\). There
\(\mathcal A=0\), \(v=K e_\theta\), and
\(\pi=-\int_r^\infty K^2\dd\rho/\rho\).
The heat-profile formula gives
\[
 |\partial_\tau^jK|\le C_jr^{-1-2h-2j},\qquad
 |\partial_\tau^j\pi|
 \le\frac{C_j}{2+4h+2j}r^{-2-4h-2j}.
\]
Their further spatial derivatives are bounded at \(r\ge r_0/2\).
Temporal cutoff derivatives occur only where
\(\tau_0/2\le\tau\le\tau_0\), which gives \(q\ge\tau_0/2\).
Every transition thus supplies finite jets for
\eqref{r:eq:localized-jet-bound}.

Near the proposed singular point both cutoffs equal one. The source is
then the local residual itself. Theorem~3.1(iii) supplies its all-orders
flatness on \(X\le X_{\rm ext}\); beyond that region the residual is
exactly zero. Combining this input with the transition bounds proves
\begin{equation}
 \sup_{x\in\R^3,\,0\le t<1}|\partial_x^\alpha\partial_t^j f(x,t)|<\infty
 \qquad\hbox{for every finite }\alpha,j.
 \label{r:eq:localized-source-jets}
\end{equation}
One additional time derivative in this bound makes each preceding jet
uniformly Cauchy as \(t\uparrow1\). Uniform spatial derivative limits
give compatible smooth terminal fields with fixed compact support.
Let \(K\subset\mathbb R^3\) be a fixed compact set containing these supports.
The source's finite Sobolev coordinates can now be bounded explicitly.
Put
\[
 M_{N,M}=\max_{0\le j\le N,\ |\alpha|\le M}
 \sup_{x,t<1}|\partial_x^\alpha\partial_t^j f(x,t)|.
\]
Expanding the integer Fourier weight \((1+|\xi|^2)^M\), whose
nonnegative multinomial coefficients sum to \(4^M\), gives
\begin{equation}
 \sum_{j=0}^N\norm{\partial_t^jf}_{L^2(0,1;H^M)}^2
 \le (N+1)4^M|K|M_{N,M}^2<\infty.
 \label{r:eq:localized-source-rectangle}
\end{equation}
Thus the actual localized source belongs to every finite mixed Sobolev
coordinate under the stated local inputs. The next construction extends
its terminal fields through time one.

Localization thus preserves source smoothness under the interior
hypotheses of \cite[Theorem~3.1]{openai-paper}. The local residual
estimate \cite[(3.4)]{openai-paper}, based on the finite-stage estimate
\cite[(9.18)]{openai-paper}, is an input to this conclusion; the
calculation above bounds the additional terms produced by the cutoffs.

\subsection{Continuing the terminal force jets}


Write \(F_j\in C_c^\infty(\R^3;\R^3)\), with
\(\operatorname{supp}F_j\subset K\), for the terminal field obtained from
\(\partial_t^j f(\cdot,t)\) as \(t\uparrow1\). The source vanishes near
\(t=0\), and all its left terminal jets have support in one fixed compact
set \(K\). These jets can be continued through time one without imposing a
bound common to all derivative orders.

Choose \(\chi_0\in C_c^\infty(\R)\), equal to one near zero and zero for
arguments at least one. For \(\sigma=t-1\ge0\), define
\begin{equation}
 \widetilde f_{\rm OA}(x,1+\sigma)=\sum_{j=0}^\infty
 \chi_0(b_j\sigma)\frac{\sigma^j}{j!}F_j(x).
 \label{r:eq:force-borel}
\end{equation}
For \(0\le t<1\), set \(\widetilde f_{\rm OA}(x,t)=f(x,t)\), the full
localized momentum residual. Formula \eqref{r:eq:force-borel} is the
extension used in \cite[(10.11)]{openai-paper}. Its full product rule is
\[
 \partial_\sigma^m
 \left(\chi_0(b_j\sigma)\frac{\sigma^j}{j!}\right)
 =\sum_{\ell=\max(0,m-j)}^m
 \binom m\ell b_j^\ell\chi_0^{(\ell)}(b_j\sigma)
 \frac{\sigma^{j-m+\ell}}{(j-m+\ell)!}.
\]
On a summand's support, \(\sigma\le b_j^{-1}\), so every product
\(b_j^\ell\sigma^{j-m+\ell}\) is at most \(b_j^{m-j}\). Writing
\[
 C_{j,m}=\sum_{\ell=\max(0,m-j)}^m
 \binom m\ell\frac{\norm{\chi_0^{(\ell)}}_\infty}{(j-m+\ell)!},
\]
we obtain
\begin{equation}
 \left\|\partial_x^\alpha\partial_\sigma^m
 \left[\chi_0(b_j\sigma)\frac{\sigma^j}{j!}F_j\right]\right\|_\infty
 \le C_{j,m}b_j^{m-j}\norm{F_j}_{C^{|\alpha|}}.
 \label{r:eq:force-borel-bound}
\end{equation}
For each \(j\ge1\), choose \(b_j>b_{j-1}\), with \(b_0=1\), so that
this bound is at most \(2^{-j}\) for all
\(|\alpha|+m\le\lfloor j/2\rfloor\). There are finitely many constraints at
that step, and \(m-j<0\) in each one. Such a choice exists for every finite
collection of norms of \(F_j\).

For fixed \(\alpha,m\), the differentiated tail beginning at
\(j\ge\max\{1,2(|\alpha|+m)\}\) is uniformly summable. At
\(\sigma=0\), the \(m\)-th temporal derivative of the series receives a
contribution only from \(j=m\), and that contribution is \(F_m\). All left
and right jets agree. Each term retains support in \(K\) and vanishes for
\(\sigma\ge1\). Uniform differentiated convergence gives smoothness at the
joining time and at the outer support boundary. Hence
\begin{equation}
 \widetilde f_{\rm OA}\in
 C_c^\infty(\R^3\times(0,\infty);\R^3)
 \subset C^\infty([0,\infty);\mathcal S(\R^3;\R^3)),
 \label{r:eq:oa-force-class}
\end{equation}
and its restriction to \(t<1\) is exactly the localized momentum residual.
For the quantitative integrals below, write
\begin{equation}
 f_{\rm OA}:=\widetilde f_{\rm OA}\big|_{[0,1)}=f.
 \label{r:eq:oa-force-restriction}
\end{equation}
Its support lies in \(K\times[0,2]\), and it retains the prescribed initial
interval of rest. A different smooth continuation beyond time one has the
same restriction on every \([0,S]\), \(S<1\), so strict-subinterval
uniqueness gives the same comparison below.
